En esta etapa de verificación, se busca garantizar que los resultados dados por
el sistema sean confiables en comparación con los datos reales.

Teniendo en cuenta que en el sistema se está trabajando con historiales de consumo mensual, es decir, series de tiempo, se decidió implementar una estrategia de validación que respetara el orden cronológico de la información, en la que los datos más antiguos siempre son los datos de entrenamiento y los más recientes se usan como datos de prueba, los cuales no se le muestran al sistema durante el entrenamiento, ya que son estos los que posteriormente tendrá que predecir.

Es por esto que se decidió implementar un sistema de validación cruzada para los
datos, en el que se repite el proceso de entrenamiento y evaluación sobre
distintas ventanas de tiempo sucesivas, para verificar que el funcionamiento del
sistema sea estable y consistente a lo largo del tiempo, y que no dependa de las
particularidades de un periodo específico para dar buenos resultados. Si al
finalizar la validación cruzada se observa que el sistema tuvo un desempeño
similar en todos los periodos evaluados, se puede concluir que ese modelo es
confiable; sin embargo, si ocurre lo contrario y solo tuvo un desempeño decente
en un periodo específico, dicho modelo no es confiable, y solo tuvo buenos
resultados en ese periodo por casualidad o por condiciones particulares del
mismo.

Para cada ventana de tiempo de la validación cruzada se calculan las métricas de MAE, RMSE, Mediana del Error, R\textsuperscript{2}, MAPE y WAPE, comparando el desempeño de los tres modelos candidatos (Media Móvil, Regresión Lineal y XGBoost) entre si.

Adicionalmente, también se evalúa la sección en la que se combinan los resultados de predicción de los modelos XGBoost y Media Móvil, en el que se estima si el promedio de ambos modelos reduce el error en comparación a usar ambos modelos por separado. Esto permite determinar si la decisión de implementar un sistema con tres modelos de predicción, y escoger el de mejor desempeño para cada producto, aporta valor en comparación con simplemente usar un solo modelo para todos los productos.
    

En cada evaluación varía el tamaño de prueba (función \texttt{tamano\_test\_backtest}), si se cuenta con un historial de consumo con una cantidad de datos más amplia, se puede usar un periodo de prueba más grande. Esto también es conveniente a la hora de evaluar la prueba de  Diebold-Mariano correspondiente a la selección del mejor modelo, ya que esta prueba ayuda a determinar si la selección de dicho modelo se dio porque de verdad es el modelo con mejores resultados o si la diferencia que tiene con el segundo mejor modelo podría tratarse solo de ruido.

En general, la estrategia de selección del modelo ganador aplica parte de la estrategia de validación, ya que solo se puede considerar un modelo como ganador para un producto si existe suficiente evidencia estadística para afirmar que su ventaja en comparación con los demas modelos es real.